\documentclass[10pt,conference]{IEEEtran}

% ============================================================
% Packages
% ============================================================
\usepackage[T1]{fontenc}
\usepackage[utf8]{inputenc}
\usepackage{times}
\usepackage{microtype}
\usepackage{graphicx}
\usepackage{booktabs}
\usepackage{array}
\usepackage{tabularx}
\usepackage{enumitem}
\usepackage{xcolor}
\usepackage{hyperref}
\usepackage{listings}
\usepackage{amsmath}
\usepackage{amssymb}

\hypersetup{
    colorlinks=true,
    linkcolor=blue,
    citecolor=blue,
    urlcolor=blue
}

\setlist[itemize]{leftmargin=*, itemsep=2pt, topsep=2pt}
\setlist[enumerate]{leftmargin=*, itemsep=2pt, topsep=2pt}
\lstset{
  basicstyle=\ttfamily\footnotesize,
  breaklines=true,
  frame=single,
  showstringspaces=false,
  columns=fullflexible,
  keepspaces=true
}

% ============================================================

% ============================================================
% Title
% ============================================================
\title{
A Lightweight Retrieval-Augmented\\
Shakespeare Assistant:\\
Domain Adaptation with a Small Language Model
}

\author{
\IEEEauthorblockN{
Group Number: \emph{Insert Group Number}
}
\IEEEauthorblockA{
CSCI433/933 Machine Learning Algorithms and Applications\\
Student Names and Student Numbers:\\
Cooper Kawari Menget (9179379), Bishal Bhatta (9990938), Raúl Palomo Mazo (1016374), Ping-Hsien Lo (1447865)\\
}
}

\begin{document}
\maketitle

% ============================================================
% Page Limit Notice
% ============================================================
% \begin{center}
% \fbox{
% \begin{minipage}{0.94\linewidth}
% \textbf{Main Report Page Limit:} The main report must not exceed \textbf{10 pages}, in the style of IEEE conference submissions. 
% The 10-page (absolute) limit applies; from the abstract through the conclusion.
% References, appendices, full evaluation tables, extended prompts, and the GenAI usage log are excluded from the 10-page limit.
% Marks are awarded for clarity, technical depth, evidence, and quality of reasoning, not for length.
% \end{minipage}
% }
% \end{center}

% ============================================================
% Abstract
% ============================================================
\begin{abstract}
We present a Shakespeare-aware question-answering system built around a
small, deployable retrieval-augmented generation (RAG) pipeline. Three
plays (\emph{Hamlet}, \emph{Macbeth}, and \emph{Romeo and Juliet}) are
indexed at the scene level using
\texttt{sentence-transformers/all-MiniLM-L6-v2} embeddings;
\texttt{phi3:3.8b} is called locally via Ollama to
compose beginner-friendly answers conditioned on the retrieved scenes.
A deterministic extractive composer is retained as a fallback when
Ollama is not running, and a stylised composer produces short
Shakespearean-style verse.
The system runs end-to-end on a standard laptop without GPU
acceleration or external API access. Across fifteen evaluation
questions (five instructor-provided, ten group-designed), the RAG
system outperformed a prompt-only baseline on every assessed
criterion: correctness $4.17$ vs $2.67$, grounding $5.00$ vs $2.43$,
retrieval relevance $4.71$ vs N/A, usefulness $5.00$ vs $2.87$, and
style quality $3.67$ vs $3.00$ (1--5 rubric), giving a composite
score of $4.63$ vs $2.66$. The most informative gains were in
grounding and usefulness, both of which the prompt-only baseline
cannot deliver. We discuss three failure modes -- weak lexical match,
phi3 paraphrasing scene citations rather than emitting them verbatim,
and graceful out-of-domain refusal -- and present an end-to-end
working system that runs on a standard laptop with no GPU.
\end{abstract}

% ============================================================
\section{Introduction and Problem Formulation}
\label{sec:introduction}

% \textbf{Relevant marks: Problem formulation and task understanding (10 marks).}

Large language models have driven recent progress in natural language
processing, but their cost, opacity, and external dependencies make
them poorly suited to organisations that need bounded, auditable,
locally deployed text systems. Small language models (SLMs), combined
with retrieval-augmented generation (RAG), offer an alternative: a
purpose-built pipeline that is cheap to run, privacy-preserving, and
easier to reason about. This assignment frames Shakespearean drama as
a controlled domain-adaptation task in which an SLM is adapted to a
narrow corpus through retrieval and prompt design rather than
parameter updates.

\textbf{Problem statement.} We build a Shakespeare-aware question
answering system over \emph{Hamlet}, \emph{Macbeth}, and
\emph{Romeo and Juliet}. The target user has little or no prior
exposure to Shakespeare and asks plain-English questions; the system
must answer them in plain English while grounding every substantive
claim in retrieved evidence. Following the assignment specification
the system supports four interaction types: \emph{concept
explanation}, \emph{contextual question answering}, \emph{evidence
retrieval}, and \emph{stylised generation}. Shakespearean text makes
this non-trivial: its lexicon and morphology (\emph{thou}, \emph{hath},
\emph{ere}) diverge from modern English, motivation often depends on
dramatic conventions and prior narrative state, and beginners need
explanation rather than paraphrase.

\textbf{Our approach.} We adopt a small-model RAG design that is
fully reproducible on a standard laptop with no GPU and no external
API. The corpus is chunked at the scene level and indexed using
\verb|sentence-transformers/all-MiniLM-L6-v2| (22M parameters); top-3
scenes are retrieved by cosine similarity and passed to phi3:3.8b
(3.8B parameters, 4-bit quantised) via local Ollama for generation.
A deterministic extractive composer provides a graceful fallback
when Ollama is unavailable, and a low-similarity threshold refuses
out-of-domain queries before they reach the generator.

\textbf{Contributions.} We report three findings.
\textit{(i)} On 15 evaluation questions (five instructor-provided,
ten group-designed), the RAG system improves composite score from
2.66 to 4.63 over a prompt-only baseline, with grounding and
usefulness rising from 2.43/2.87 to 5.00/5.00.
\textit{(ii)} A two-line citation instruction in the system prompt
lifts grounding from 3.29 to 4.57 without affecting correctness or
retrieval, demonstrating that prompt engineering can be measured
independently of retrieval quality.
\textit{(iii)} A 0.15 cosine-similarity threshold reliably refuses an
out-of-domain probe (cosine $0.085$) while remaining permissive on
all in-scope questions (cosine $> 0.55$).

\textbf{Report structure.} Section~\ref{sec:dataset} describes the
dataset preparation and chunking strategy. Section~\ref{sec:methodology}
covers the baseline and RAG pipeline.
Section~\ref{sec:system} describes the system implementation and
how to reproduce the results. Section~\ref{sec:evaluation-design}
defines the evaluation question set and the scoring rubric.
Section~\ref{sec:results} presents the results and per-type
breakdown, and discusses four failure modes. Sections~\ref{sec:genai}
and~\ref{sec:conclusion} cover responsible GenAI use and conclusions.

% ============================================================
\section{Dataset and Preprocessing}
\label{sec:dataset}

An existing structured Shakespeare dialogue set was used for the project that included all three required plays (Hamlet, Macbeth and Romeo and Juliet).

The set was developed for RAG and modeled the plays as hierarchy of scenes of utterances, metadata and context: Each scene comprised a number of utterances with speaker labels, dialogue text, scene id, and other metadata. Additional context was given through the \texttt{scene\_summary} and \texttt{keywords}.

The scene summaries provided a modern-English description of what occurred in the scene.The keywords provided a list of prominent subjects, topics, or characters to aid in retrieval. Normalizing: During the normalization stage of preprocessing, the corpus was normalized to provide more consistent retrieval representations and to unneeded meta-data.The original data set contained both sourceid and utteranceid which were discovered to be always equal, and So there was no need to have them distinct, so only one was supplied as the retrieval chunk id. In the same way, both speaker and speaker original data were supplied but as all speaker original data was only used in the sense of upper/lower case and abbreviated versions (e.g. Ber and bernardo) the normalised speaker data was used. Whitespace and formatting issues were processed with text normalisation using regular-expression substitution.

Multiple whitespace and newline characters were compressed to a single space to improve embedding uniformity and avoid retrieval noise.Records with null or badly encoded fields were discarded. The non character theatrical markers and structural artifacts that were not useful for retrieval were further filtered out.

These included labels such as {FLOURISH} or {THUNDER} or {ALARUM} or {MUSIC} or location headings such as {ELSINORE} or {VERONA}.

\begin{verbatim}
	[STAGE_DIRECTION] Enter Ghost.
\end{verbatim}

This allowed the embedding model to distinguish between spoken dialogue and narrative stage actions while preserving important contextual information.

The preprocessing pipeline therefore transformed the dataset from a raw literary transcription into a retrieval-oriented representation suitable for semantic search and evidence-grounded question answering.


\subsection{Dataset schema}

After preprocessing, each retrieval record retained only the metadata required for retrieval, traceability, and evaluation. The processed schema preserved contextual information including play, act, scene, speaker identity, summaries, and source identifiers.

An example processed record is shown below:
\begin{quote}
	\small
	\begin{ttfamily}
		\begin{flushleft}
			\{
			\newline
			"play": "Macbeth",
			\newline
			"act": 1,
			\newline
			"scene": 3,
			\newline
			"speaker": "MACBETH",
			\newline
			"scene\_summary": "Macbeth and Banquo encounter the witches.",
			\newline
			"keywords": ["witches", "prophecy", "Banquo"],
			\newline
			"text": "So foul and fair a day I have not seen.",
			\newline
			"chunk\_id": "macbeth\_1\_3\_0001"
			\newline
			\}
		\end{flushleft}
	\end{ttfamily}
\end{quote}
Each processed retrieval chunk preserved:
\begin{itemize}
	\item \texttt{play}: source play title;
	\item \texttt{act} and \texttt{scene}: narrative location within the play;
	\item \texttt{speaker}: normalised speaker label;
	\item \texttt{scene\_summary}: modern-English contextual summary;
	\item \texttt{keywords}: semantic topic hints;
	\item \texttt{text}: dialogue or stage direction content;
	\item \texttt{chunk\_id}: traceable source identifier.
\end{itemize}

The inclusion of summaries and keywords provided supplementary semantic information that improved embedding quality and retrieval recall, particularly for beginner-oriented questions where users may not use Shakespearean wording directly.
\subsection{Chunking strategy}


The system primarily used scene-level chunking. In this strategy, all utterances belonging to a single scene were combined into one retrieval chunk. Scene-level chunking was selected because Shakespearean dialogue is highly contextual and meaning often depends on surrounding interactions between multiple speakers. Questions such as ``Why is Horatio afraid?'' or ``Why does Macbeth hesitate?'' require narrative context beyond a single line of dialogue.

Each scene chunk included:
\begin{itemize}
	\item speaker dialogue;
	\item stage directions;
	\item scene summaries;
	\item keywords;
	\item play, act, and scene metadata.
\end{itemize}

This allowed retrieval to preserve conversational flow, emotional progression, and narrative continuity. The scene summaries and keywords were prepended to the chunk text before embedding generation so that the vector representation captured both the original Early Modern English dialogue and a simplified semantic explanation of the scene.

An alternative utterance-level chunking strategy was also implemented for comparison purposes. In utterance chunking, each speaker turn became an individual retrieval chunk. This approach produced smaller and more precise retrieval units that were effective for direct quote lookups or speaker identification tasks. However, utterance-level chunks frequently lacked sufficient narrative context to answer explanatory ``why'' or ``how'' questions.

The trade-offs between the two strategies are summarised below:

\begin{table}[h]
	\centering
    \caption{Comparison of chunking strategies}
    \label{tab:chunking}
	\small
	\begin{tabularx}{\linewidth}{|p{2cm}|X|X|}
		\hline
		\textbf{Strategy} & \textbf{Advantages} & \textbf{Disadvantages} \\
		\hline
		Scene-level chunking &
		Preserves context and narrative flow; stronger for semantic reasoning &
		Larger chunks may introduce retrieval noise \\
		\hline
		Utterance-level chunking &
		High precision for quote retrieval; smaller embeddings &
		Frequently loses surrounding context \\
		\hline
	\end{tabularx}
\end{table}
Other chunking strategies, such as fixed-size token chunking, recursive chunking, and semantic topic-based chunking, were considered but not implemented in the final system. Fixed-size chunking can arbitrarily divide dialogue and stage directions, potentially separating important conversational context across chunk boundaries. Recursive and semantic chunking approaches may preserve semantic coherence more effectively, but they introduce additional preprocessing complexity and computational overhead that was unnecessary for the relatively well-structured Shakespeare dataset.

Since the plays are already naturally divided into acts and scenes, scene-level structural chunking provided a simpler and more interpretable retrieval representation while still preserving narrative continuity. For the scope of this assignment, scene-level chunking represented an effective balance between retrieval quality, contextual preservation, implementation simplicity, and explainability. Utterance-level chunking was retained as a secondary strategy primarily for diagnostic comparison and direct quote retrieval tasks.

% ============================================================
\section{Methodology}
\label{sec:methodology}

% \textbf{Relevant marks: Model adaptation and RAG methodology (20 marks).}

\subsection{Baseline system}

The baseline, \verb|PromptOnlyBaseline| (\verb|src/baseline.py|),
returns a fixed, generic answer drawn from a small lookup of
play-level priors (e.g., a one-sentence summary of \emph{Macbeth}).
It does \emph{not} call the retriever and therefore cannot cite a
scene, act, or speaker. The baseline is intentionally weak: its
purpose is to expose the contribution of retrieval rather than to
compete on absolute quality. A second variant,
\verb|KeywordRetrievalBaseline|, uses Jaccard overlap over scene
summaries; it is included in the repository for completeness but is
not the primary baseline reported here. We chose the prompt-only
variant as the primary baseline because the assignment specification
explicitly lists ``prompt-only generation'' as an acceptable baseline.

\subsection{RAG-based system}

The RAG system (\verb|src/rag_chatbot.py|) has five components.

\textbf{Embedder / retriever.} Implemented in
\verb|src/retrieval.py|. The class \verb|EmbeddingRetriever| has a
single public API (\verb|build_index|, \verb|retrieve|) and two
interchangeable backends:
\textit{(a)} \verb|sentence-transformers/all-MiniLM-L6-v2| dense
embeddings (22M parameters, 384-d, CPU-friendly), used when the
package is installed;
\textit{(b)} a scikit-learn TF--IDF backend (unigrams + bigrams,
English stop-word removal, \verb|sublinear_tf=True|, L2 normalisation)
used otherwise. The active backend is exposed via
\verb|retriever.backend_name| and recorded in
\verb|results/evaluation_summary.json|. All reported results in this
report use the dense MiniLM backend; the summary JSON confirms
\texttt{backend\_name} as \texttt{sentence-transformers/all-MiniLM-L6-v2}.
The TF--IDF backend is retained as a zero-dependency fallback and is
the simplest implementation listed as acceptable in the assignment
specification.

\textbf{Top-$k$.} We use $k = 3$. With $73$ scene chunks, this returns
roughly one scene per play on average; a beginner-friendly answer
typically needs only the highest-scoring scene as evidence.

\textbf{Prompt structure.} \verb|src/rag_chatbot.py| builds a prompt
that injects the retrieved evidence between a system prompt
(\verb|prompts/system_prompt.txt|, instructing the model to use only
retrieved context, indicate uncertainty, and avoid invented details)
and the user query. The prompt is materialised even when the
extractive composer is used so that the system can be transparently
re-wired to a hosted SLM later (see \verb|HuggingFaceCausalGenerator|
in \verb|src/generation.py|).

\textbf{Generator: phi3:3.8b via Ollama (default).} The default
generator is Microsoft's \texttt{phi3:3.8b} (3.8B parameters,
instruction-tuned, 4-bit quantised) called over the local Ollama HTTP
API at \verb|http://localhost:11434|. The model receives the system
prompt, the top-$k$ retrieved scenes, and the user question (see
``Prompt structure'' above) and returns a beginner-friendly answer
constrained by the prompt. We picked phi3:3.8b because it is one of
the smallest competent instruction-tuned models that runs on a
standard laptop without a GPU, satisfies the assignment's
``small, deployable, resource-constrained'' framing, and is fully
local (no data leaves the host).

\textbf{Fallback generator: extractive composer.} If Ollama is not
running, the system silently falls back to a deterministic extractive
composer (\verb|generation.ExtractiveComposer|) that constructs each
answer from four ingredients drawn solely from the retrieved scenes:

\begin{enumerate}
    \item an opener cued by the query type
        (\emph{contextual\_qa}, \emph{concept\_explanation},
        \emph{stylised\_generation});
    \item the modern-English scene summary(ies) of the top scenes;
    \item the relevant keyword set;
    \item a single best-quoted sentence picked from the highest-scoring
        scene by lexical overlap with the query (capped at
        \verb|MAX_QUOTE_WORDS|=60).
\end{enumerate}

The composer never emits text that is not in (or trivially derived
from) the retrieved chunks, so the system has a strong grounding
guarantee by construction. When the top similarity falls below
\verb|min_score|=0.025 the composer prints an explicit
``the retrieved passages do not appear to address this question with
high confidence'' message together with the closest match, so weak
retrieval is surfaced rather than hidden.

\textbf{Generator: stylised composer.}
\verb|generation.StylisedComposer| produces a $\leq$150-word
Shakespearean-style verse stitched together from short quoted seeds
plus a small set of templated transitions, with a light register
shift (\texttt{you} $\to$ \texttt{thou}, \texttt{before} $\to$
\texttt{ere}, \ldots). Every stylised output is prefixed with
``\texttt{[STYLISED CREATIVE OUTPUT -- not textual evidence; based on
the retrieved scenes.]}'' as required by the specification.

\subsection{Model choice and justification}

We deliberately avoided an external hosted API and a large local
generative model. The assignment's pedagogical focus is on
``\emph{purpose-built systems that can operate under constraints involving
cost, privacy, latency, maintainability, and domain specificity}''.
A deterministic extractive composer aligns with all five constraints:
zero inference cost, no data leaves the host, sub-second latency,
trivial maintainability, and guaranteed domain specificity (the system
cannot mention plays outside the corpus). The retrieval backend
(MiniLM, 22M parameters, 384-dimensional, CPU-only) is similarly
small for an embedding model; the TF--IDF fallback is smaller still
(a few hundred kilobytes of fitted weights and a sparse-vector cosine
product at inference).

We did, however, design the code so that an SLM call can be slotted
in without rewriting the pipeline. The class
\verb|HuggingFaceCausalGenerator| in \verb|src/generation.py| wraps a
small instruction-tuned model (default
\verb|google/flan-t5-small|) behind the same interface. We did not
evaluate that path in this report because the marking is over the
default configuration we can demonstrably run, but the section
``Possible extensions'' returns to this point.

% ============================================================
\section{System Design and Implementation}
\label{sec:system}

% \textbf{Relevant marks: System implementation and usability (15 marks).}

\subsection{System pipeline}

The end-to-end flow per user query is:

\begin{enumerate}
    \item user query received via CLI (\verb|rag_chatbot.py|) or
        evaluation harness (\verb|evaluate.py|);
    \item the query is embedded with the same backend used at
        index time;
    \item top-$k$ scene chunks are retrieved by cosine similarity;
    \item question type is inferred by a small rule classifier
        (\verb|classify_question|);
    \item the appropriate composer runs
        (\emph{stylised} or \emph{extractive});
    \item the retrieved evidence is displayed alongside the generated
        answer (CLI) and recorded in the evaluation log (CSV + JSONL).
\end{enumerate}

\subsection{Repository structure and how to run}

\begin{lstlisting}[basicstyle=\ttfamily\scriptsize,frame=none,xleftmargin=0pt]
assessment 2/
  data/processed/      # JSON + JSONL processed plays
  prompts/             # system_prompt.txt, stylised_prompt.txt
  results/             # instructor + group questions,
                       # evaluation_results.csv / .jsonl /
                       # evaluation_summary.json
  src/
    config.py          # paths, defaults
    data_loader.py     # JSON -> flat utterance records
    chunking.py        # scene / utterance chunkers
    retrieval.py       # SBERT (preferred) or TF-IDF
    generation.py      # Extractive + Stylised composers
    baseline.py        # PromptOnly / KeywordRetrieval baselines
    rag_chatbot.py     # full RAG chatbot, CLI entry point
    build_index.py     # sanity-check script
    evaluate.py        # eval harness writing the result files
  report/              # this report
  requirements.txt
\end{lstlisting}

To reproduce the evaluation from scratch:

\begin{lstlisting}[basicstyle=\ttfamily\scriptsize,frame=none,xleftmargin=0pt]
pip install -r requirements.txt
cd src
python build_index.py     # smoke test
python evaluate.py        # writes results/*.{csv,jsonl,json}
python rag_chatbot.py     # interactive CLI
\end{lstlisting}

\subsection{Caching and reproducibility}

The retrieval index is rebuilt from the JSON corpus on each run.
With \verb|sentence-transformers/all-MiniLM-L6-v2| the 73 scene chunks
embed in roughly 10--15\,s on a laptop CPU; with the TF--IDF fallback
the index is rebuilt in under a second. In both cases the index is
deterministic given the input corpus, so we do not persist it
between runs --- the on-disk artefacts that matter for
reproducibility are the JSON corpus (unchanged) and the source code
(versioned). phi3:3.8b is called via Ollama with
\texttt{temperature}$=0.2$, introducing minor non-determinism in
generation; observed score variation across re-runs was negligible
(see Sec.~\ref{sec:evaluation-design}).

\subsection{Implementation limitations}

\begin{itemize}
    \item phi3:3.8b is a free-form generative model. We mitigate
        hallucination via the retrieval-first prompt structure, but
        cannot fully eliminate it -- this is the central
        accuracy/grounding trade-off discussed in
        Sec.~\ref{sec:failure-analysis}. The deterministic extractive
        composer is retained as a fallback for stricter grounding.
    \item Dense embeddings (MiniLM) sometimes retrieve the
        correct play but the wrong scene when the query phrasing
        does not surface in the target scene's wording --- see
        Failure 1 in Sec.~\ref{sec:failure-analysis}. The TF--IDF
        fallback shares this paraphrase-gap limitation in stronger
        form.
    \item Out-of-domain detection relies on a single cosine
        similarity threshold (\verb|min_retrieval_score|$=0.15$).
        This caught the banana probe (G6, top similarity $0.085$)
        but a malicious or borderline-relevant query that scores
        slightly above the threshold would still be forwarded to
        phi3, which may then answer from world knowledge.
\end{itemize}

% ============================================================
\section{Evaluation Design}
\label{sec:evaluation-design}

% \textbf{Relevant marks: Evaluation, comparison, and error analysis (20 marks).}

\subsection{Question set}

The evaluation corpus comprises \textbf{15 questions}: five
instructor-provided (Q1--Q5) and ten group-designed (G1--G10), stored
in \verb|results/instructor_questions.json| and
\verb|results/group_questions.json| respectively.
Table~\ref{tab:question-types} shows the distribution by type.

\begin{table}[h]
\centering
\caption{Evaluation question distribution by type.}
\label{tab:question-types}
\small
\begin{tabular}{lcc}
\toprule
Type & $n$ & Questions \\
\midrule
\texttt{contextual\_qa}       & 5 & Q1, Q4, G1, G3, G4 \\
\texttt{concept\_explanation} & 4 & Q2, Q3, G2, G5     \\
\texttt{stylised\_generation} & 3 & Q5, G7, G10         \\
\texttt{evidence\_retrieval}  & 2 & G8, G9              \\
\texttt{out\_of\_domain}      & 1 & G6                  \\
\bottomrule
\end{tabular}
\end{table}

The group questions are designed to test aspects of the system not
exercised by Q1--Q5:

\begin{itemize}
    \item \textbf{G1, G2} -- secondary-character depth (Lady Macbeth;
        the witches) whose key scenes span multiple acts.
    \item \textbf{G3, G4} -- multi-scene reasoning across
        \emph{Romeo and Juliet}.
    \item \textbf{G5} -- the explicit ``in plain English'' modifier,
        testing whether the system honours a beginner-friendliness cue.
    \item \textbf{G6} -- a deliberate out-of-domain probe
        (\emph{``What's a banana?''}). A scope-aware system should
        refuse rather than answer from general world knowledge.
    \item \textbf{G7} -- a second stylised-generation request
        (Macbeth monologue), paired with Q5 (Juliet), to isolate
        style quality from character choice.
    \item \textbf{G8, G9} -- scene-specific evidence retrieval:
        \emph{``Find the scene where Hamlet speaks to Yorick's skull''}
        (expected: Act~5, Scene~1) and \emph{``Show me the passage
        where Lady Macbeth tries to wash imaginary blood from her
        hands''} (expected: Act~5, Scene~1). These stress-test the
        retriever's \emph{scene-level} precision, not just play-level
        relevance.
    \item \textbf{G10} -- a third stylised-generation request
        (Hamlet grief soliloquy), bringing the stylised sample to
        three for within-type comparison.
\end{itemize}

\subsection{Scoring criteria}

Five criteria are scored on a 1--5 scale and implemented
deterministically in \verb|src/evaluate.py|. The automated scorer
emits values in $\{1, 3, 5\}$; intermediate scores (2, 4) require
manual annotation and may be recorded in the \texttt{comments} column
of the CSV per the rubric in \verb|docs/scoring_rubric.md|.

\begin{itemize}
    \item \textbf{Correctness} (1, 3, 5): fraction of
        \texttt{expected\_focus} content words present in the answer:
        $\geq$0.5 $\to$ 5; $>$0 $\to$ 3; 0 $\to$ 1. Not scored for
        \texttt{stylised\_generation} (N/A). For
        \texttt{out\_of\_domain}, a correctly detected refusal
        scores 5; an answer scores 1.

    \item \textbf{Grounding} (1, 3, 5): 5 if the answer contains
        an explicit ``Act~$i$, Scene~$j$'' citation; 3 if it names a
        play; 1 otherwise. Not scored for \texttt{out\_of\_domain}.

    \item \textbf{Retrieval relevance} (1, 3, 5): for
        \texttt{contextual\_qa}, \texttt{concept\_explanation}, and
        \texttt{stylised\_generation}, 5 if the top retrieved chunk is
        from the expected play, 3 if any retrieved chunk is, 1
        otherwise. For \texttt{evidence\_retrieval} questions the
        scorer requires an exact act-and-scene match for 5; returning
        the right play but the wrong scene scores 3. This distinction
        is implemented via \verb|_parse_expected_scene()| in
        \verb|evaluate.py|, which extracts the expected act/scene from
        the \texttt{expected\_focus} free text and passes it to
        \verb|score_retrieval()|. Baseline rows and
        \texttt{out\_of\_domain} rows are marked N/A.

    \item \textbf{Usefulness} (1, 3, 5): 5 if the answer is
        $\geq$25 informative words \emph{and} carries an explicit
        citation (grounding $=$ 5); 3 if either condition fails; 1 if
        the answer matches a refusal phrase. For
        \texttt{out\_of\_domain}, a refusal scores 5 and an answer
        scores 1.

    \item \textbf{Style quality} (1, 3, 5): for
        \texttt{stylised\_generation} questions only -- 5 if the
        answer contains $\geq$3 archaic markers (\texttt{thou},
        \texttt{thy}, \texttt{ere}, \texttt{o'er}, \texttt{hark},
        \texttt{thine}, \texttt{doth}, \texttt{didst}); 3 if
        $\geq$1; 1 otherwise.
\end{itemize}

A weighted composite score is computed per question as:
\[
    c = \frac{\displaystyle\sum_{k \in \mathcal{A}} w_k \cdot s_k}
             {\displaystyle\sum_{k \in \mathcal{A}} w_k}
\]
where $\mathcal{A}$ is the set of applicable criteria for that
question type, with weights $w_\text{corr} = 0.30$,
$w_\text{ground} = 0.25$, $w_\text{retr} = 0.20$,
$w_\text{use} = 0.20$, $w_\text{style} = 0.05$.
N/A criteria are excluded and the remaining weights are renormalised.

A \textbf{potential-hallucination flag} is raised when the RAG system
cites a play or scene (grounding $\geq$ 3) but the retrieved chunks
are from the wrong play (retrieval relevance $=$ 1), indicating a
citation not supported by the actual retrieved evidence.
No potential hallucinations were detected in the final evaluation run.

\subsection{Two-stage scoring methodology}

The five criteria above are applied in two complementary stages that
together constitute the full evaluation.

\textbf{Stage 1 -- automated scoring (\texttt{evaluate.py}).}
The harness scores every response programmatically using surface-level
heuristics: word-overlap ratios, regular expressions for citation
patterns, and keyword presence. This produces reproducible, auditable
scores that can be regenerated from scratch by running
\texttt{python evaluate.py}. The fundamental constraint of this
approach is that a heuristic can only make reliable three-way
distinctions: clearly wrong, partial, clearly right. So the
automated scores are restricted to $\{1, 3, 5\}$. Attempting to emit
2 or 4 automatically would introduce arbitrary thresholds and a false
impression of precision.

\textbf{Stage 2 -- human review (\texttt{docs/scoring\_rubric.md}).}
The rubric defines the full 1--5 scale for human evaluators who read
each response in context. A human can apply genuine judgment, for
example, ``this answer cites the correct play and roughly the right
scene but omits one key character'' warrants a 4 on grounding, not a
3, in a way that no regex can. The rubric provides a consistent
reference point so that all team members score the same row the same
way. Manual adjustments and their rationale are recorded in the
\texttt{comments} column of \verb|evaluation_results.csv|; final
summary means should reflect the manually adjusted values.

\textbf{Why this design.} Separating automated from manual scoring is
standard practice in NLP system evaluation~\cite{lewis2020rag}. The
automated stage provides coverage and reproducibility across all 30
rows at negligible cost; the manual stage adds the precision and
semantic understanding that heuristics cannot provide. The two
artefacts are therefore complementary: \verb|evaluate.py| is not a
replacement for the rubric, nor is the rubric a critique of the
automated scorer. Each does what the other cannot.

\subsection{Evaluation procedure}

The evaluation harness (\verb|src/evaluate.py|) runs both systems on
every question in a single pass and writes three artefacts:

\begin{itemize}
    \item \verb|evaluation_results.csv| -- 30 rows
        (15 questions $\times$ 2 systems), one row per
        question--system pair;
    \item \verb|evaluation_results.jsonl| -- same rows with full
        chunk metadata (rank, cosine similarity, play/act/scene);
    \item \verb|evaluation_summary.json| -- criterion means per
        system and per question type, plus composite scores,
        generator provenance, and the potential-hallucination list.
\end{itemize}

\textbf{Generator transparency.} Every CSV row carries a
\texttt{generator\_path} column recording which generator produced
the response: \texttt{ollama/phi3:3.8b} (normal path),
\texttt{extractive\_fallback} or \texttt{stylised\_fallback} (Ollama
unavailable), \texttt{low\_similarity\_refusal} (retrieval score
below threshold), or \texttt{prompt\_only\_baseline}. The summary
JSON records the union of generator paths used across the run.
In the reported results, \texttt{generators\_used} contains
\texttt{["low\_similarity\_refusal", "ollama/phi3:3.8b",
"prompt\_only\_baseline"]}, confirming that phi3 generated every
in-scope RAG response and that the refusal path fired exactly once
(G6).

\textbf{Out-of-domain threshold.}
A \texttt{min\_retrieval\_score} of 0.15 is applied before calling
phi3. If the top retrieved chunk scores below this threshold the
system returns a canned refusal rather than forwarding the query to
phi3, which would otherwise answer from general world knowledge. The
banana probe (G6) returned a top cosine similarity of 0.085;
all in-scope Shakespeare questions scored above 0.55.

\subsection{Limitations of the scoring approach}

\begin{itemize}
    \item \textbf{Coarse auto-scorer.} The automated heuristics
        produce only $\{1, 3, 5\}$. Intermediate values (2, 4) require
        manual annotation. The scorer is intentionally simple so it is
        transparent and fully reproducible, but may under- or
        over-score edge cases.

    \item \textbf{Correctness is token-overlap only.} The metric
        counts \texttt{expected\_focus} content words in the answer.
        It rewards lexical matches but cannot assess semantic
        accuracy: a wrong answer that repeats the expected terms would
        score 5.

    \item \textbf{Retrieval is scene-level for
        \texttt{evidence\_retrieval}, play-level elsewhere.}
        For G8 and G9 the scorer correctly penalises wrong-scene
        retrieval (score 3 rather than 5); for all other question
        types only the play is checked. The evidence-retrieval
        retrieval-relevance mean of 3.0 reflects a genuine system
        limitation: \texttt{all-MiniLM-L6-v2} retrieves the correct
        play reliably but, in both cases, returns an earlier scene
        rather than Act~5, Scene~1. This is discussed further in
        Sec.~\ref{sec:failure-analysis}.

    \item \textbf{phi3 temperature.} phi3:3.8b is called with
        \texttt{temperature}$=0.2$, introducing minor
        non-determinism. In practice we observed negligible score
        variation across re-runs. All reported results are from a
        single canonical run with \texttt{generators\_used}
        confirmed in the summary JSON.

    \item \textbf{Small question set.} With 15 questions, per-type
        means rest on 1--5 items and should be interpreted with
        caution. The \texttt{evidence\_retrieval} composite of 4.579
        is the average of exactly two questions.
\end{itemize}

% ============================================================
\section{Results and Analysis}
\label{sec:results}

\textbf{Relevant marks: Evaluation, comparison, and error analysis (20 marks).}

\subsection{Baseline versus RAG}

Table~\ref{tab:evaluation-summary} reports mean criterion scores
across all 15 questions (applicable rows only; N/A criteria excluded
from each mean). The RAG system outperforms the prompt-only baseline
on every assessed criterion, with a composite score of $4.626$ versus
$2.663$.

\begin{table}[h]
\centering
\caption{Mean scores (1--5), 15 questions. Retrieval relevance is N/A
for the baseline (no retriever). Composite weights:
$w_\text{corr}=0.30$, $w_\text{grnd}=0.25$, $w_\text{ret}=0.20$,
$w_\text{use}=0.20$, $w_\text{sty}=0.05$.}
\label{tab:evaluation-summary}
\small
\begin{tabularx}{\linewidth}{lXXXXXX}
\toprule
System & Corr. & Grnd. & Retr. & Use. & Style & Comp. \\
\midrule
Baseline         & 2.67 & 2.43 & N/A  & 2.87 & 3.00 & 2.663 \\
RAG (phi3+MiniLM)& \textbf{4.17} & \textbf{5.00} & \textbf{4.71} & \textbf{5.00} & \textbf{3.67} & \textbf{4.626} \\
\midrule
$\Delta$         & $+$1.50 & $+$2.57 & --- & $+$2.13 & $+$0.67 & $+$1.963 \\
\bottomrule
\end{tabularx}
\end{table}

The largest absolute gains are in \textbf{grounding} ($+2.57$) and
\textbf{usefulness} ($+2.13$). Grounding reached the maximum mean of
$5.0$: the system prompt's explicit citation instruction (``Always
cite the play, act, and scene number'') caused phi3 to emit structured
``Act~$i$, Scene~$j$'' references in virtually every factual response,
satisfying the rubric's top-level criterion.
Usefulness follows grounding mechanically, the heuristic scores 5
only when both length ($\geq$25 words) and an explicit citation are
present, so the grounding improvement propagated directly.
\textbf{Correctness} rises by $+1.50$: phi3 synthesises an actual
answer from the retrieved context rather than copying a scene
summary verbatim, so the expected-focus content words appear more
reliably. \textbf{Style quality} improves by $+0.67$, though phi3's
verse is inconsistent (discussed in Sec.~\ref{sec:failure-analysis}).

\subsection{Per-type breakdown}

Table~\ref{tab:by-type} breaks down RAG performance by question type.
Retrieval relevance is $5.0$ for all types where the expected play can
be inferred, except \texttt{evidence\_retrieval} ($3.0$), which
requires scene-level precision (Sec.~\ref{sec:failure-analysis}).
\texttt{out\_of\_domain} achieves a perfect composite of $5.0$: the
similarity-threshold guard correctly refused the out-of-scope query
(G6, cosine similarity $0.085$), avoiding the general-knowledge
hallucination that the unguarded system produced.

\begin{table}[h]
\centering
\caption{RAG mean scores by question type (1--5). N/A: criterion not
applicable to type. $n$ = number of questions.}
\label{tab:by-type}
\small
\begin{tabularx}{\linewidth}{lXXXXXX}
\toprule
Type & $n$ & Corr & Grnd & Retr & Use & Comp \\
\midrule
\texttt{concept\_expl.}   & 4 & 4.50 & 5.00 & 5.00 & 5.00 & 4.842 \\
\texttt{contextual\_qa}   & 5 & 3.40 & 5.00 & 5.00 & 5.00 & 4.495 \\
\texttt{evidence\_retr.}  & 2 & 5.00 & 5.00 & 3.00 & 5.00 & 4.579 \\
\texttt{out\_of\_domain}  & 1 & 5.00 & N/A  & N/A  & 5.00 & 5.000 \\
\texttt{stylised\_gen.}   & 3 & N/A  & 5.00 & 5.00 & 5.00 & 4.905 \\
\bottomrule
\end{tabularx}
\end{table}

The weakest per-type correctness is \texttt{contextual\_qa} ($3.40$):
phi3 answers causal ``why?'' questions coherently but does not always
reproduce the specific tokens named in \texttt{expected\_focus},
reducing the overlap score. This reflects a genuine limitation of
the token-overlap correctness metric rather than factual error in
most cases.

\subsection{Qualitative examples}

\textbf{Q1 (``Why does Macbeth kill Duncan?'' -- success).}
The RAG system retrieves \emph{Macbeth}, Act~2, Scene~2 (similarity
$0.758$) and Act~1, Scene~4 ($0.730$). phi3 synthesises a
beginner-friendly answer, quotes Lady Macbeth directly
(``That which hath made them drunk hath made me bold''), and cites
Act~2, Scene~2 explicitly. Scores: correctness~5, grounding~5,
retrieval~5, usefulness~5. The baseline returns a generic
one-sentence play summary with no citation.

\textbf{G6 (``What's a banana?'' -- out-of-domain refusal).}
Top retrieved chunk similarity $0.085$, below the
\texttt{min\_retrieval\_score} threshold of $0.15$. The system
returns: \emph{``This question is outside the scope of the Shakespeare
corpus (Hamlet, Macbeth, Romeo and Juliet). No relevant passages were
retrieved.''} Scores: correctness~5, usefulness~5. The baseline
answers with a generic disclaimer but does not refuse.

\textbf{G7 (``Write a Macbeth monologue'' -- stylised success).}
phi3 generates a Shakespearean-style monologue grounded in the
witches' prophecy (Act~1, Scene~3), citing it explicitly and including
archaic markers (\emph{quoth}, \emph{forsooth}, \emph{mine},
\emph{thee}, \emph{doth}). Prefixed with the mandatory
``[STYLISED CREATIVE OUTPUT]'' disclosure. Scores: grounding~5,
retrieval~5, usefulness~5, style~5.

\textbf{G8 (``Find the Yorick skull scene'' -- retrieval gap).}
The retriever returns Act~2, Scene~2 and Act~4, Scene~2 -- correct
play, wrong scenes. phi3 nevertheless answers correctly (citing
Act~5, Scene~1) by drawing on its training knowledge. Correctness~5,
grounding~5, but retrieval relevance~3. This is a notable behaviour:
phi3 compensates for retrieval failure with parametric knowledge,
which is useful here but is a form of ungrounded generation -- the
cited evidence was not in the retrieved context.

\subsection{Failure analysis}
\label{sec:failure-analysis}

\textbf{Failure 1: scene-level retrieval gap (G8, G9).}
Both evidence-retrieval questions target Act~5, Scene~1 -- the
graveyard scene in \emph{Hamlet} and the sleepwalking scene in
\emph{Macbeth}. In both cases MiniLM retrieves the correct play but
an earlier scene (Act~2 and Act~4 in \emph{Hamlet}; Act~1 and Act~2
in \emph{Macbeth}). The late scenes score poorly on cosine similarity
with a query phrased in modern English: the corpus summaries for
Act~5 do not repeat the lexical cues in the question
(``skull'', ``wash'', ``blood''), while earlier scenes containing
those themes score higher. A position-weighted re-ranker or
scene-level keyword index would directly address this; it is the
clearest extension for future work.

\textbf{Failure 2: phi3 compensates with training knowledge.}
When retrieval fails (G8, G9), phi3 draws on its own training rather
than the retrieved context. For G8 this produces a correct answer;
for G9, phi3 generates a plausible-looking but fabricated passage
attributed to Lady Macbeth. Neither case triggered the
potential-hallucination flag (which requires a grounding/retrieval
mismatch), because phi3 emits citations from the retrieved chunk
headers regardless of whether the cited scene was actually retrieved.
This is a fundamental limitation of the retrieval-first prompt
design: a fluent generative model with strong prior knowledge can
bypass the grounding constraint.

\textbf{Failure 3: contextual\_qa correctness ceiling (mean 3.40).}
``Why?'' questions require phi3 to synthesise multiple causal factors.
For Q4 (``Why does Hamlet delay taking revenge?'') phi3 generates a
coherent answer about Hamlet's philosophical indecision but cites
Act~4, Scene~5 (a non-standard scene number) and does not reproduce
the expected-focus tokens \emph{uncertainty}, \emph{moral hesitation},
\emph{proof}. The response is substantively correct but scores
correctness~1 under the token-overlap heuristic. This is the most
visible gap between the automated correctness score and human judgment.

\textbf{Failure 4: style quality inconsistency (G10 scores 1).}
For G10 (``Hamlet expresses grief in Shakespearean style'') phi3
returns a short verbatim quote
(``Alas, poor Yorick! I knew him'') with no archaic register markers
beyond the quote itself. The response scores style~1 because it
contains fewer than the required three markers. phi3's stylised
output quality varies across questions: G7 (Macbeth monologue) scores
style~5, while G10 and Q5 score~3 and~5 respectively. The
inconsistency likely reflects phi3's sensitivity to the exact phrasing
of the creative prompt.

\textbf{Issues identified and resolved during development.}
Two failure modes were found and fixed before the final evaluation
run and are reported here for completeness.
\emph{(a) Out-of-domain hallucination:} without the retrieval
threshold, phi3 returned a factually correct description of a banana
for G6, entirely ignoring the Shakespeare corpus. Fixed by adding the
\texttt{min\_retrieval\_score}~$=0.15$ guard in
\verb|src/rag_chatbot.py|.
\emph{(b) Grounding paraphrasing:} phi3 cited play names but not
act/scene numbers, scoring grounding~$\approx$3.3 on average. Fixed
by adding an explicit citation instruction to
\verb|prompts/system_prompt.txt|, raising mean grounding to $5.0$.

% ============================================================
\section{Responsible Use of Generative AI}
\label{sec:genai}

% \textbf{Relevant marks: Responsible and transparent GenAI usage (10 marks).}

Generative AI assistance was used for code scaffolding, naming, and
debugging. Every artefact in the repository -- code, prompts, the
group-designed evaluation questions, and the text of this report --
was reviewed, modified, and signed off by the group. The full GenAI
usage log is included in Appendix~\ref{app:genai-log}.

In summary: AI was used to (i) suggest the dual-backend retrieval
architecture (MiniLM as primary with a TF--IDF fallback behind one
interface), (ii) draft an initial version of the evaluation rubric
which we subsequently revised to discrete $\{1, 3, 5\}$ scoring with
human-review escape hatches for 2 and 4, and (iii) sketch the
archaic register-shift heuristics for the stylised composer. In
every case the suggestion was checked against the assignment
specification, executed end-to-end, and modified where the output
diverged from what the group wanted (see the verification column of
Table~\ref{tab:genai-log}). We did not paste AI output unverified
into the report or codebase.

% ============================================================
\section{Conclusion}
\label{sec:conclusion}

We built a small, fully deployable Shakespeare RAG system that runs
on a standard laptop with no external API and no GPU. The system
satisfies all four interaction types required by the specification:
concept explanation, contextual question answering, evidence
retrieval, and stylised generation under a 150-word cap. Across 15
evaluation questions the RAG system improves the composite score
from $2.66$ to $4.63$ versus a prompt-only baseline, with grounding
rising from $2.43$ to $5.00$ and usefulness from $2.87$ to $5.00$
on a 1--5 rubric.

The most consequential design decisions were
\textit{(i)} scene-level chunking augmented with the
instructor-provided scene summary and keywords, which gave the
retriever enough modern-English signal to match plain-English
queries against Early Modern English dialogue;
\textit{(ii)} an explicit citation instruction in the system prompt,
which lifted grounding by $1.7$ points without affecting correctness
or retrieval; and
\textit{(iii)} a cosine-similarity threshold (\verb|min_retrieval_score|$=0.15$)
that refuses out-of-domain queries before phi3 can hallucinate from
its training data.

The chief residual limitations are
\textit{(a)} scene-level retrieval precision: MiniLM reliably finds
the correct play but, for the two evidence-retrieval probes, returned
earlier scenes rather than the iconic Act~5, Scene~1 scenes we
expected;
\textit{(b)} phi3 occasionally compensates for retrieval failure with
parametric training knowledge, producing a fluent but ungrounded
answer that the current potential-hallucination flag does not catch;
and
\textit{(c)} the token-overlap correctness scorer under-credits
synthesised answers that are substantively right but do not reuse
the expected-focus tokens, evident in the \texttt{contextual\_qa}
type mean of $3.40$.

Given more time we would (a) add a hybrid BM25 + cross-encoder
re-ranker, or a question-type-aware retrieval boost, to close the
scene-level gap; (b) tighten the hallucination check to compare
phi3's cited act/scene against the actually retrieved chunks rather
than just the play; and (c) build a small human-annotated gold set
so the rubric can credit semantically correct but lexically
mismatched answers.

% ============================================================
% References
% ============================================================
\begin{thebibliography}{1}

\bibitem{reimers2019sbert}
N.~Reimers and I.~Gurevych, ``Sentence-BERT: Sentence embeddings using
Siamese BERT-networks,'' in \emph{Proc. EMNLP-IJCNLP}, 2019.

\bibitem{lewis2020rag}
P.~Lewis et~al., ``Retrieval-augmented generation for
knowledge-intensive NLP tasks,'' in \emph{Proc. NeurIPS}, 2020.

\bibitem{robertson2009bm25}
S.~Robertson and H.~Zaragoza, ``The probabilistic relevance framework:
BM25 and beyond,'' \emph{Found. Trends Inf. Retr.}, vol.~3, no.~4,
pp.~333--389, 2009.

\bibitem{scikitlearn}
F.~Pedregosa et~al., ``Scikit-learn: Machine learning in Python,''
\emph{J. Mach. Learn. Res.}, vol.~12, pp.~2825--2830, 2011.

\end{thebibliography}

% ============================================================
% Appendices
% ============================================================
\appendices

\section{Full Evaluation Table}
\label{app:evaluation}

The complete evaluation table is generated by \verb|src/evaluate.py|
and stored as \verb|results/evaluation_results.csv| (15 questions
$\times$ 2 systems = 30 rows). Each row records the question, the
retrieved passages, the generated response, all five criterion scores,
the \texttt{generator\_path}, and a \texttt{potential\_hallucination}
flag.

Table~\ref{tab:full-eval-rag} shows the RAG system scores per
question. ``--'' marks criteria not applicable to the question type
(correctness is N/A for stylised generation; style is N/A for
non-stylised; grounding and retrieval are N/A for the
\texttt{low\_similarity\_refusal} path used on G6).

\begin{table}[h]
\centering
\caption{Per-question RAG scores (1--5). Means: Corr 4.17, Grnd 5.00, Retr 4.71, Use 5.00, Style 3.67. Composite 4.626.}
\label{tab:full-eval-rag}
\small
\begin{tabularx}{\linewidth}{lXccccc}
\toprule
ID & Type (abbrev.) & Corr & Grnd & Retr & Use & Styl \\
\midrule
Q1 & contextual\_qa     & 5 & 5 & 5 & 5 & -- \\
Q2 & concept\_expl.     & 3 & 5 & 5 & 5 & -- \\
Q3 & concept\_expl.     & 5 & 5 & 5 & 5 & -- \\
Q4 & contextual\_qa     & 1 & 5 & 5 & 5 & -- \\
Q5 & stylised\_gen.     & -- & 5 & 5 & 5 & 5 \\
G1 & contextual\_qa     & 3 & 5 & 5 & 5 & -- \\
G2 & concept\_expl.     & 5 & 5 & 5 & 5 & -- \\
G3 & contextual\_qa     & 5 & 5 & 5 & 5 & -- \\
G4 & contextual\_qa     & 3 & 5 & 5 & 5 & -- \\
G5 & concept\_expl.     & 5 & 5 & 5 & 5 & -- \\
G6 & out\_of\_domain    & 5 & --& --& 5 & -- \\
G7 & stylised\_gen.     & -- & 5 & 5 & 5 & 5 \\
G8 & evidence\_retr.    & 5 & 5 & 3 & 5 & -- \\
G9 & evidence\_retr.    & 5 & 5 & 3 & 5 & -- \\
G10& stylised\_gen.     & -- & 5 & 5 & 5 & 1 \\
\bottomrule
\end{tabularx}
\end{table}

The per-row CSV remains the authoritative record; the summary JSON
\verb|results/evaluation_summary.json| also records the retrieval
backend as \texttt{sentence-transformers/all-MiniLM-L6-v2}, the
indexed chunk count as 73, and \texttt{generators\_used} as
\texttt{["low\_similarity\_refusal", "ollama/phi3:3.8b",
"prompt\_only\_baseline"]}.

The baseline scores are omitted from this table for space but are
available in the CSV. In every row the baseline scores 1 on
grounding-related criteria, with retrieval marked N/A by construction
(no retriever).

\section{Representative Prompts}
\label{app:prompts}

System prompt (\verb|prompts/system_prompt.txt|):

\begin{lstlisting}[basicstyle=\ttfamily\scriptsize,frame=single,xleftmargin=0pt]
You are a Shakespeare-aware assistant.
Use the retrieved context to answer the user's
question.
Your answer must be beginner-friendly.
Always cite the play, act, and scene number
when you reference retrieved text (e.g., "In
Hamlet, Act 3, Scene 1..."). Include at least
one such citation in every factual answer.
The retrieved text is from public domain works;
you may quote it directly.
If the retrieved context is insufficient, say
so clearly.
Do not invent unsupported details.
If asked to produce a Shakespearean-style
response, keep it under 150 words and make
clear that it is creative output, not textual
evidence.
\end{lstlisting}

Stylised prompt (\verb|prompts/stylised_prompt.txt|, used to brief
the SLM in stylised-generation mode):

\begin{lstlisting}[basicstyle=\ttfamily\scriptsize,frame=single,xleftmargin=0pt]
You are a Shakespeare-aware assistant producing
a short stylised response.
Rules:
- The response must be at most 150 words.
- Adopt early-modern English register ("thou",
  "thy", "ere", "o'er", "doth").
- The response is creative output and must be
  labelled as such; it must not be presented as
  textual evidence.
- The response should remain interpretable to a
  beginner reader.
- Use the retrieved scenes only as inspiration
  for the themes touched on (ambition, love,
  prophecy, etc.); do not quote them as if they
  were the user's question.
\end{lstlisting}

At inference, the full RAG prompt concatenates the system prompt,
the top-$k$ retrieved chunks (each annotated with
``[Context $i$ $|$ similarity=$s$]'' headers), and the user query
under an ``Answer:'' label. The complete prompt is materialised in
\verb|src/rag_chatbot.py| and recorded for every evaluation row in
\verb|results/evaluation_results.jsonl|.

\section{Additional Output Examples}
\label{app:outputs}

Three complete RAG outputs are reproduced below, illustrating one
successful concept question, one out-of-domain refusal, and one
stylised-generation case.

\textbf{Q3 (concept explanation): ``What is the conflict between
the Montagues and the Capulets?''}

\begin{lstlisting}[basicstyle=\ttfamily\scriptsize,frame=single,xleftmargin=0pt,breaklines=true]
In Romeo and Juliet, Act 1, Scene 1, we learn
about a longstanding feud. Verona's citizens
are caught in this bitter rivalry as Sampson
of the house of Montague says to Gregory of
the same family, "But thou art not quickly
moved to strike." This indicates that even
among friends and relatives within their own
families (Montagues), there is tension due to
the feud with the Capulets.
\end{lstlisting}

Retrieved evidence: \emph{Romeo and Juliet}, Act~1, Scene~2
(similarity $0.556$) and Act~1, Scene~1 ($0.554$). Scores:
correctness~5, grounding~5, retrieval~5, usefulness~5. Note that
phi3 successfully cites the play, act, and scene, and quotes a line
from the retrieved Act~1, Scene~1 chunk verbatim. The reading of
``Gregory of the same family'' as fellow Montagues is correct.

\textbf{G6 (out-of-domain refusal): ``What's a banana?''}

\begin{lstlisting}[basicstyle=\ttfamily\scriptsize,frame=single,xleftmargin=0pt,breaklines=true]
This question is outside the scope of the
Shakespeare corpus (Hamlet, Macbeth, Romeo and
Juliet). No relevant passages were retrieved.
\end{lstlisting}

Top retrieved chunk: \emph{Romeo and Juliet}, Act~2, Scene~3,
similarity $0.085$. Below the \texttt{min\_retrieval\_score}
threshold of $0.15$; \texttt{generator\_path} is recorded as
\texttt{low\_similarity\_refusal}, so phi3 is never invoked.
Scores: correctness~5 (correctly refused), usefulness~5.

\textbf{G7 (stylised generation): ``Write a short Shakespearean
monologue from Macbeth after he hears the witches' prophecy.''}

\begin{lstlisting}[basicstyle=\ttfamily\scriptsize,frame=single,xleftmargin=0pt,breaklines=true]
[STYLISED CREATIVE OUTPUT -- not textual
evidence; generated by phi3:3.8b grounded on
the retrieved scenes.]

"Hark! Three sisters of shadows and deceit
have spoken unto me their cryptic words,
foretelling my ascent to kingship by a lineage
not born but conjured in ambition's crucible.
'Give me,' quoth I with eager heart, 'aroint
thee,' they say! Yet herein lies the rub: Is
it fate or folly that drives mine hand towards
this throne? Forsooth, these prophecies do
stir my soul to action; yet in their depths
doth lurk doubt. Shall ambition's fruit be
sweet when sown with treachery's seed? I am
but a man of mortal flesh and blood -- a tyrant
not ordained by heavenly decree, nor born from
the womb that cradles life. What is this power
they promise me if it comes at such perilous
cost to my soul's integrity?"

(This monologue reflects Macbeth's internal
conflict upon hearing the prophecy in Act 1,
Scene 3 of "Macbeth.")
\end{lstlisting}

Retrieved evidence: \emph{Macbeth}, Act~1, Scene~3
(similarity $0.728$) and Act~5, Scene~3 ($0.710$). Scores:
grounding~5, retrieval~5, usefulness~5, style~5 (the response
contains \emph{Hark}, \emph{quoth}, \emph{mine}, \emph{Forsooth},
\emph{doth}, \emph{thy} -- well above the 3-marker threshold).

The full set of generated outputs and the corresponding retrieved
chunks (with scene IDs, ranks, and similarity scores) is available
in \verb|results/evaluation_results.jsonl|.

\section{GenAI Usage Log}
\label{app:genai-log}

This appendix documents how generative AI tools (primarily Claude and
ChatGPT) were used during the project. Per the assignment specification,
each row records the task, the useful AI suggestion, and the
verification or modification applied by the group before adoption.

\begin{table*}[h]
\centering
\caption{GenAI usage log.}
\label{tab:genai-log}
\footnotesize
\begin{tabularx}{\textwidth}{p{2.2cm} p{3.0cm} p{3.4cm} X}
\toprule
Tool & Prompt or Task & Useful Output & Verification / Modification \\
\midrule
Claude
& Asked how to fix the \texttt{minted} package error during LaTeX compilation.
& Suggested switching from \texttt{minted} to \texttt{listings}, which ships with LaTeX, with install-free setup instructions.
& Verified by re-compiling; \texttt{pdflatex} succeeded after the swap. The \texttt{listings} configuration is now in the preamble of this report. \\
\midrule
Claude / ChatGPT
& ``How can I support both sentence-transformers and a TF--IDF fallback behind one interface?''
& Proposed a thin \texttt{EmbeddingRetriever} abstraction with a common \texttt{encode(texts)} method on both backends.
& Adopted. We kept the public API unchanged but added a \texttt{backend\_name} attribute and surfaced it in \verb|evaluation_summary.json| so every reported run records which backend was used. \\
\midrule
Claude / ChatGPT
& ``Draft a 1--5 scoring rubric for grounding, retrieval relevance, correctness, usefulness, and style.''
& Initial draft with continuous 5-point scales.
& Tightened to discrete \{1, 3, 5\} for the automated scorer, with 2 and 4 reserved for human review (documented in \verb|docs/scoring_rubric.md|). Added the curly-quote normaliser and a refusal-phrase list after observing false negatives. \\
\midrule
Claude / ChatGPT
& ``Suggest archaic-register substitutions for a stylised generator.''
& Pairs such as \texttt{you}~$\to$~\texttt{thou}, \texttt{before}~$\to$~\texttt{ere}, \texttt{over}~$\to$~\texttt{o'er}.
& Adopted conservatively. Rejected substitutions that produced ungrammatical English (e.g.\ \texttt{has}~$\to$~\texttt{hath} without subject-agreement logic); kept only the substitutions that round-trip safely. \\
\midrule
ChatGPT
& Discuss preprocessing and chunking strategies for the Shakespeare RAG system.
& Suggested metadata normalisation, scene-level chunking justification, and preprocessing improvements.
& Suggestions were reviewed and selectively adapted to fit the implemented system and assignment requirements. \\
\midrule
ChatGPT
& Assist with LaTeX formatting and debugging in TeXstudio.
& Provided fixes for table formatting, overflow issues, and compilation errors.
& All formatting changes were manually tested and adjusted in the final report. \\
\midrule
Claude / ChatGPT
& ``Help me parse expected act/scene from the \texttt{expected\_focus} free text so the retrieval scorer can do scene-level matching.''
& Suggested a regex covering ``Act 5 Scene 1'', ``Act 5, Scene 1'', Roman numerals, and ``5.1''.
& Adopted. We added an assertion in \verb|run_evaluation()| that fails fast on \texttt{evidence\_retrieval} questions whose \texttt{expected\_focus} cannot be parsed, so question-authoring bugs surface immediately rather than silently scoring 5. \\
\midrule
Claude
& Report layout review on IEEEtran template (column widths, table placement, caption position).
& Pointed out that \texttt{\textbackslash textwidth} in two-column overflows, that \texttt{verbatim} does not break long lines, and that IEEE convention places table captions above tables.
& Applied to every table and code block in this report; all wide tables now use \texttt{\textbackslash linewidth} or \texttt{table*} as appropriate. \\
\bottomrule
\end{tabularx}
\end{table*}

In every case the AI suggestion was treated as a starting point rather
than a finished artefact: code was executed end-to-end, prompts were
iterated, and any output that did not behave as the group intended was
modified or discarded. We did not paste AI output unverified into the
report or codebase.

\end{document}
